\section{Protected low-leakage phase: low temperature as suppressed exit (extended interface module)}
\label{app:protected_low_leakage_phase}

\paragraph{Scope and status.}
\AuditTag This appendix records a disciplined way to talk about ``very low temperature'' in the tick-trap language without assuming a specific microscopic condensed-matter mechanism.
The phrase ``crystal/superconducting state'' is treated only as an explanatory label.
The theorem-level folding core is not affected.

\subsection{Primary signature (A): suppressed leakage}
\label{subsec:pll_primary_signature}

\begin{definition}[Protected low-leakage phase (audit definition)]
\label{def:protected_low_leakage_phase}
Fix a target $X$, a coarse-graining that exposes a finite exit-channel set $\mathcal C$, and a finite candidate family of leakage kernels (Definition~\ref{def:survival_kernel_family}).
We say that $X$ is in a \emph{protected low-leakage phase} (relative to this audit setting) if the CAP-selected leakage rate proxy satisfies
\[
\Gamma_{\mathrm{CAP}}(X)\ll \Gamma_{\mathrm{ref}},
\]
for a declared reference scale $\Gamma_{\mathrm{ref}}$ appropriate to the comparison class, and the conclusion is robust under the declared finite counterfactual family (e.g.\ alternative coarse-grainings or kernel-family variants).
\end{definition}

\paragraph{Interpretation.}
\AuditTag The operational content is: exits are extremely rare; dwell times are extremely long.
This is the direct tick-first translation of ``low effective temperature'' as ``weak emission/leakage'' at the interface.
No claim is made that the underlying mechanism is literally a BCS condensate or a crystalline lattice.

\subsection{Derived diagnostic (B): near-lossless transport response}
\label{subsec:pll_derived_transport}

\paragraph{Near-elastic response.}
\AuditTag When leakage is strongly suppressed, the external response is dominated by phase accumulation rather than absorption.
In a scattering readout, this typically manifests as a nearly unitary response with small absorptive part over the relevant band.
This motivates using near-lossless transport as a \emph{diagnostic} of the low-leakage phase.
\AuditTag On such near-unitary bands, the delay proxy is also a logdet derivative of the scattering response (Lemma~\ref{lem:wigner_smith_trace_logdet}), so loss-control (S1) is simultaneously the gatekeeping condition for interpreting $\tau_{\mathrm{WS}}$ as a stable phase/determinant-based readout rather than as a lossy proxy.

\paragraph{Connection to delay.}
\AuditTag In the mass-as-latency dictionary (Section~\ref{sec:mass_latency_coordinate}), protected traps can simultaneously exhibit large Wigner--Smith delay $\tau_{\mathrm{WS}}$ (deep dwell) while still having extremely small leakage rate $\Gamma$ (rare exits).
The pair $(\tau_{\mathrm{WS}},\Gamma)$ is therefore the minimal audit signature used downstream.

\subsection{Specialization under \texorpdfstring{$P_{\mathrm{phys}}$}{Pphys}: black-hole low temperature}
\label{subsec:pll_bh_specialization}

\AuditTag Under the boundary/Planck information pack (Assumptions~\ref{ass:bh_planck_units_info}--\ref{ass:bh_schwarzschild_package}), Hawking temperature decreases with mass for Schwarzschild black holes.
In this specialization, ``low temperature'' has an external meaning via $T_H(M)$.
The role of the present appendix is to provide the protocol translation: the low-$T$ regime is modeled as a regime of strongly suppressed leakage (small $\Gamma$) with long dwell (large $\tau$), and any further claim about a specific ``superconducting'' mechanism must be introduced as an additional optional model assumption in Appendix~\ref{app:physics_consensus_inputs}.

